\documentclass{article}
\usepackage[preprint]{neurips_2023}
\usepackage[utf8]{inputenc}
\usepackage[T1]{fontenc}
\usepackage{microtype}
\usepackage[hidelinks]{hyperref}
\setcitestyle{numbers,square}
\makeatletter
\renewcommand{\@noticestring}{Research concept note. September 2026.}
\makeatother
\title{Learning the Meanings of Justice}
\author{A proposed collaboration in language representation learning and transitional justice}
\begin{document}
\maketitle

An ongoing Columbia research program offers an opportunity to develop language models that capture how people express justice, harm and reparation. The proposed collaboration would build on the work of \textbf{Valentina Velásquez}, M.S. candidate at the School of Professional Studies, and \textbf{Peter Dixon}, Associate Professor of Practice and Director of Columbia's M.S. in Negotiation and Conflict Resolution, and Senior Fellow at the Center for Justice.

\textbf{The technical question.} Can domain-informed human supervision improve how encoders represent specific aspects of meaning across new accounts? An encoder maps a passage to a vector used for comparison, retrieval and clustering. Conditional semantic similarity asks how that comparison should change with the question being asked. Two narratives might share an initiating event while describing different consequences. C-STS formalizes conditional similarity~\citep{deshpande2023}; Tu et al. generate aspect-specific answers before comparison~\citep{tu2024}; CASE learns a compact transformation of conditional embeddings~\citep{zhang2026}. These provide concrete baselines for testing the value of expert supervision.

\textbf{Exploratory analysis.} I built and ran a reproducible pilot with 24 constructed narrative triplets. Each test required selecting the passage that shared an anchor's initiating event, then the passage that shared its consequence. Generating an aspect-specific answer before embedding increased Qwen's joint accuracy---the share of triplets with both selections matching the intended answers---from 42\% (10/24) to 75\% (18/24). AI judges consistently endorsed both intended answers on seven triplets, motivating expert annotation to clarify the distinctions and establish reliable evaluation labels.

\textbf{A shared research opportunity.} Columbia's School of Professional Studies (SPS) team assembled \textbf{243 hearing videos spanning more than 1,500 hours} and analyzed 40 videos in its initial pilot. Its summaries report 2,145 coded justice claims, 25 coding categories and a 30-segment sample reviewed by experts. These team-reported findings span JEP Cases 01 and 03~\citep{team}. Its workflow combines AI extraction, human validation, codebook development and expert review. The Casanare extension examines reparation, restoration, sanction and harm. Combining this foundation and the team's Spanish-language expertise with JEP expert annotations could produce an \textbf{expert-annotated benchmark}. LLMs could surface informative comparisons and disagreements, while experts supply judgments and supporting evidence. We would test how this supervision improves conditional encoders, then use retrieval and clustering to explore recurring meanings across the corpus.

\textbf{Next steps.} The SPS team is seeking technical oversight for its AI-assisted research. An initial collaboration with Columbia Computer Science could help shape the annotation protocol, select model baselines, and establish a reproducible evaluation. Faculty sponsorship would guide a focused pilot using the team's corpus and expert judgments, connecting research in representation learning with the substantive work of interpreting victims' accounts and informing transitional justice.

\begin{thebibliography}{4}
\small
\bibitem{deshpande2023} Deshpande et al. \href{https://aclanthology.org/2023.emnlp-main.345/}{C-STS: Conditional Semantic Textual Similarity}. In \emph{Proceedings of EMNLP}, 2023.
\bibitem{tu2024} Tu et al. \href{https://aclanthology.org/2024.acl-long.64/}{Linguistically Conditioned Semantic Textual Similarity}. In \emph{Proceedings of ACL}, 2024.
\bibitem{zhang2026} Zhang et al. \href{https://aclanthology.org/2026.eacl-long.231/}{CASE---Condition-Aware Sentence Embeddings for Conditional Semantic Textual Similarity Measurement}. In \emph{Proceedings of EACL}, 2026.
\bibitem{team} SPS research team. \emph{Contexto del Proyecto}; \emph{Script LSA presentation}; proposed workflow outline. Unpublished project materials supplied by the team.
\end{thebibliography}
\end{document}
